\documentclass[11pt,a4paper]{article}
\usepackage[utf8]{inputenc}
\usepackage[T1]{fontenc}
\usepackage{amsmath,graphicx,hyperref}
\usepackage[margin=2.5cm]{geometry}
\title{Un espacio local para escribir}
\author{Proyecto de ejemplo}
\date{\today}
\begin{document}
\maketitle
\tableofcontents
\section{Introducci\'on}\label{sec:introduccion}
Este proyecto permite probar el editor sin documentos personales.
Haz doble clic sobre este texto en el PDF para localizarlo en su fuente.

El contenido se guarda en la carpeta vinculada y las versiones se conservan
en \texttt{.latex-local/}. La referencia \cite{Demo2026} es sint\'etica y
solo sirve para mostrar el flujo bibliogr\'afico.

\section{Una tabla de ejemplo}
\begin{table}[h]
\centering
\begin{tabular}{ll}
\hline
Herramienta & Uso \\
\hline
Historial & Comparar versiones \\
Referencias & Incorporar BibTeX \\
SyncTeX & Navegar desde el PDF \\
\hline
\end{tabular}
\caption{Funciones para trabajar en un documento.}\label{tab:funciones}
\end{table}
La Tabla~\ref{tab:funciones} resume algunas herramientas.
\bibliographystyle{plain}
\bibliography{references}
\end{document}
